% GENERATED FILE. Edit canonical lessons, then run npm run generate:books.
% canonical-source-hash: fnv1a64:3a3450e9cb560360
% canonical-lessons: UR-C169-bansuri, UR-C169-nach, UR-C169-numaish, UR-C169-mela, UR-C169-jashn

\chapter{Flute, Dance, Exhibition, Fair, Celebration}
\label{ch:ur-flute-dance-exhibition-fair-celebration}
\hlchaptermodality{169}

\begin{chapteropening}
\textbf{By the end of this chapter:} \emph{I can say a flute, a dance, an exhibition, a fair and a celebration.}

Complete the last lesson of chapter 169: I can say a flute, a dance, an exhibition, a fair and a celebration.
\end{chapteropening}

\section[\texorpdfstring{\ur{بانسری} (bāṅsurī)}{بانسری (bāṅsurī)}]{\ur{بانسری} (bāṅsurī) — a flute}

\label{lesson:UR-C169-bansuri}



\begin{quote}
Before the new one: say the Urdu for cricket, then the Urdu for football.
\end{quote}

\subsection*{You'll want to know: \ur{بانسری}}
\textbf{\ur{بانسری}} — \emph{bāṅsurī} — \textquotedblleft{}a flute\textquotedblright{}.

\begin{cousinweb}
\textbf{\ur{بانسری}} (\emph{bāṅsurī}) belongs with Hindi \emph{bãsī}, \textquotedblleft{}flute\textquotedblright{}, which CDIAL lists under Sanskrit \emph{vaṃśī}, \textquotedblleft{}flute, pipe\textquotedblright{}, built on \emph{vaṃśa}, \textquotedblleft{}bamboo\textquotedblright{} (CDIAL 11180). CDIAL does not list \emph{bāṅsurī} itself.
\end{cousinweb}

\begin{grammarlens}[title={What you've built}]
One more word for this chapter.
\end{grammarlens}

\subsection*{Your turn}
\begin{itemize}
\raggedright
  \item \emph{Say it:} \emph{bāṅsurī}
  \item \emph{Say it:} \emph{bāṅsurī}, once more
  \item \emph{Say it:} say \emph{bāṅsurī}
  \item \emph{Recall:} say the Urdu for cricket, then the Urdu for football, then say \emph{bāṅsurī} again
\end{itemize}

\begin{tcolorbox}[breakable,colback=teal!4,colframe=teal!35!black,title={Before you move on}]
What does \ur{بانسری} mean? (\textquotedblleft{}A flute\textquotedblright{}.) Say it once more.
\end{tcolorbox}
\section[\texorpdfstring{\ur{ناچ} (nāch)}{ناچ (nāch)}]{\ur{ناچ} (nāch) — a dance}

\label{lesson:UR-C169-nach}



\begin{quote}
Before the new one: say the Urdu for football, then the Urdu for a flute.
\end{quote}

\subsection*{You'll want to know: \ur{ناچ}}
\textbf{\ur{ناچ}} — \emph{nāch} — \textquotedblleft{}a dance\textquotedblright{}.

\begin{cousinweb}
\textbf{\ur{ناچ}} (\emph{nāch}) is a tadbhava, a word inherited from Sanskrit and reshaped by sound change. It comes from \emph{nṛtya}, \textquotedblleft{}dancing\textquotedblright{} (CDIAL 7582). It is built on \emph{nṛt-}, \textquotedblleft{}to dance\textquotedblright{}, the root of \emph{nāchnā}.
\end{cousinweb}

\begin{grammarlens}[title={What you've built}]
One more word for this chapter.
\end{grammarlens}

\subsection*{Your turn}
\begin{itemize}
\raggedright
  \item \emph{Say it:} \emph{nāch}
  \item \emph{Say it:} \emph{nāch}, once more
  \item \emph{Say it:} say \emph{nāch}
  \item \emph{Recall:} say the Urdu for football, then the Urdu for a flute, then say \emph{nāch} again
\end{itemize}

\begin{tcolorbox}[breakable,colback=teal!4,colframe=teal!35!black,title={Before you move on}]
What does \ur{ناچ} mean? (\textquotedblleft{}A dance\textquotedblright{}.) Say it once more.
\end{tcolorbox}
\section[\texorpdfstring{\ur{نمائش} (numāish)}{نمائش (numāish)}]{\ur{نمائش} (numāish) — an exhibition}

\label{lesson:UR-C169-numaish}



\begin{quote}
Before the new one: say the Urdu for a flute, then the Urdu for a dance.
\end{quote}

\subsection*{You'll want to know: \ur{نمائش}}
\textbf{\ur{نمائش}} — \emph{numāish} — \textquotedblleft{}an exhibition\textquotedblright{}.

\begin{cousinweb}
\textbf{\ur{نمائش}} (\emph{numāish}) is a loan, usually traced to Persian \emph{numāyish}, \textquotedblleft{}display, a showing\textquotedblright{}. It is built on \emph{numūdan}, \textquotedblleft{}to show\textquotedblright{}.
\end{cousinweb}

\begin{grammarlens}[title={What you've built}]
One more word for this chapter.
\end{grammarlens}

\subsection*{Your turn}
\begin{itemize}
\raggedright
  \item \emph{Say it:} \emph{numāish}
  \item \emph{Say it:} \emph{numāish}, once more
  \item \emph{Say it:} say \emph{numāish}
  \item \emph{Recall:} say the Urdu for a flute, then the Urdu for a dance, then say \emph{numāish} again
\end{itemize}

\begin{tcolorbox}[breakable,colback=teal!4,colframe=teal!35!black,title={Before you move on}]
What does \ur{نمائش} mean? (\textquotedblleft{}An exhibition\textquotedblright{}.) Say it once more.
\end{tcolorbox}
\section[\texorpdfstring{\ur{میلہ} (melā)}{میلہ (melā)}]{\ur{میلہ} (melā) — a fair}

\label{lesson:UR-C169-mela}



\begin{quote}
Before the new one: say the Urdu for a dance, then the Urdu for an exhibition.
\end{quote}

\subsection*{You'll want to know: \ur{میلہ}}
\textbf{\ur{میلہ}} — \emph{melā} — \textquotedblleft{}a fair\textquotedblright{}.

\begin{cousinweb}
\textbf{\ur{میلہ}} (\emph{melā}) is a tadbhava, a word inherited from Sanskrit and reshaped by sound change. It comes from \emph{mela}, \textquotedblleft{}meeting\textquotedblright{} (CDIAL 10331). It is built on \emph{mil-}, \textquotedblleft{}to meet\textquotedblright{}, the root of \emph{milnā}.
\end{cousinweb}

\begin{grammarlens}[title={What you've built}]
One more word for this chapter.
\end{grammarlens}

\subsection*{Your turn}
\begin{itemize}
\raggedright
  \item \emph{Say it:} \emph{melā}
  \item \emph{Say it:} \emph{melā}, once more
  \item \emph{Say it:} say \emph{melā}
  \item \emph{Recall:} say the Urdu for a dance, then the Urdu for an exhibition, then say \emph{melā} again
\end{itemize}

\begin{tcolorbox}[breakable,colback=teal!4,colframe=teal!35!black,title={Before you move on}]
What does \ur{میلہ} mean? (\textquotedblleft{}A fair\textquotedblright{}.) Say it once more.
\end{tcolorbox}
\section[\texorpdfstring{\ur{جشن} (jashn)}{جشن (jashn)}]{\ur{جشن} (jashn) — a celebration}

\label{lesson:UR-C169-jashn}



\begin{quote}
Before the new one: say the Urdu for an exhibition, then the Urdu for a fair.
\end{quote}

\subsection*{You'll want to know: \ur{جشن}}
\textbf{\ur{جشن}} — \emph{jashn} — \textquotedblleft{}a celebration\textquotedblright{}.

\begin{cousinweb}
\textbf{\ur{جشن}} (\emph{jashn}) is a loan, usually traced to Persian \emph{jashn}, \textquotedblleft{}a feast, a festival\textquotedblright{}.
\end{cousinweb}

\begin{grammarlens}[title={What you've built}]
One more word for this chapter.
\end{grammarlens}

\subsection*{Your turn}
\begin{itemize}
\raggedright
  \item \emph{Say it:} \emph{jashn}
  \item \emph{Say it:} \emph{jashn}, once more
  \item \emph{Say it:} say \emph{jashn}
  \item \emph{Recall:} say the Urdu for an exhibition, then the Urdu for a fair, then say \emph{jashn} again
\end{itemize}

\begin{tcolorbox}[breakable,colback=teal!4,colframe=teal!35!black,title={Before you move on}]
What does \ur{جشن} mean? (\textquotedblleft{}A celebration\textquotedblright{}.) Say it once more.
\end{tcolorbox}
